\documentclass[11pt]{article}
\usepackage[utf8]{inputenc}
\usepackage[T1]{fontenc}
\usepackage{amsmath}
\usepackage{amssymb}
\usepackage{multicol}
\usepackage{geometry}
\usepackage{xcolor}
\usepackage{titlesec}

\geometry{a4paper, left=1cm, right=1cm, top=0.8cm, bottom=1.2cm}

\titleformat{\section}{\normalfont\Large\bfseries\color{blue}}{}{0em}{}
\titleformat{\subsection}{\normalfont\large\bfseries\color{red}}{}{0em}{}

\title{\textcolor{red}{Segundo Trimestre – Matemática\\Razão e Proporção}}
\author{Professor: Jefferson}
\date{}

\begin{document}

\maketitle

\begin{multicols}{2}

\section*{Questões e Gabarito Comentado}

\subsection*{Questão 1}
A razão entre dois números é $3:5$ e a soma deles é 64. Determine esses números.

\textbf{Resolução:}
Sejam os números $3x$ e $5x$.
\[
3x + 5x = 64 \Rightarrow 8x = 64 \Rightarrow x = 8
\]
Logo, os números são $24$ e $40$.

\textbf{Resposta:} 24 e 40

\subsection*{Questão 2}
A razão entre dois números é $\frac{3}{4}$. Sabendo que o menor é 45, determine o maior.

\textbf{Resolução:}
\[
\frac{3}{4} = \frac{45}{x}
\Rightarrow 3x = 180
\Rightarrow x = 60
\]

\textbf{Resposta:} 60

\subsection*{Questão 3}
Dois números estão na razão $2:7$ e a soma deles é 81. Determine o maior número.

\textbf{Resolução:}
\[
2x + 7x = 81 \Rightarrow 9x = 81 \Rightarrow x = 9
\]
Maior número:
\[
7x = 63
\]

\textbf{Resposta:} 63

\subsection*{Questão 4}
Dois números estão na razão $5:8$. Sabendo que a diferença entre eles é 27, determine o maior número.

\textbf{Resolução:}
\[
8x - 5x = 27 \Rightarrow 3x = 27 \Rightarrow x = 9
\]
Maior número:
\[
8x = 72
\]

\textbf{Resposta:} 72

\subsection*{Questão 5}
Em uma mistura, a razão entre álcool e água é $3:5$. Se a mistura possui 32 litros, quantos litros de álcool há?

\textbf{Resolução:}
Total de partes: $3 + 5 = 8$.
\[
\frac{3}{8} \times 32 = 12
\]

\textbf{Resposta:} 12 litros

\subsection*{Questão 6}
Uma receita utiliza farinha, açúcar e manteiga na razão $5:3:2$. Se foram usados 600 g de farinha, quantos gramas de açúcar foram utilizados?

\textbf{Resolução:}
\[
\frac{3}{5} = \frac{x}{600}
\Rightarrow x = 360
\]

\textbf{Resposta:} 360 g

\subsection*{Questão 7}
Um mapa utiliza a escala $1:250\,000$. Qual é a distância real correspondente a 3,2 cm no mapa?

\textbf{Resolução:}
\[
3{,}2 \times 250\,000 = 800\,000 \text{ cm}
\]
Convertendo:
\[
800\,000 \text{ cm} = 8 \text{ km}
\]

\textbf{Resposta:} 8 km

\subsection*{Questão 8}
A razão entre o número de alunos de duas turmas é $7:9$. Sabendo que a diferença entre elas é 12 alunos, determine o total de alunos.

\textbf{Resolução:}
\[
9x - 7x = 12 \Rightarrow 2x = 12 \Rightarrow x = 6
\]
Total:
\[
7x + 9x = 16x = 96
\]

\textbf{Resposta:} 96 alunos

\subsection*{Questão 9}
Um carro percorre 180 km com 12 litros de combustível. Quantos litros serão necessários para percorrer 75 km?

\textbf{Resolução:}
\[
\frac{180}{75} = \frac{12}{x}
\Rightarrow 180x = 900
\Rightarrow x = 5
\]

\textbf{Resposta:} 5 litros

\subsection*{Questão 10}
Se 3 kg de tinta pintam 45 m², quantos metros quadrados podem ser pintados com 8 kg?

\textbf{Resolução:}
\[
\frac{3}{8} = \frac{45}{x}
\Rightarrow 3x = 360
\Rightarrow x = 120
\]

\textbf{Resposta:} 120 m²

\subsection*{Questão 11}
Um trabalhador recebeu R\$ 1.260 por 18 dias de trabalho. Quanto receberá por 25 dias, mantendo a mesma proporção?

\textbf{Resolução:}
\[
\frac{18}{25} = \frac{1260}{x}
\Rightarrow 18x = 31\,500
\Rightarrow x = 1\,750
\]

\textbf{Resposta:} R\$ 1.750

\subsection*{Questão 12}
A razão entre as áreas de dois quadrados é $9:16$. Qual é a razão entre seus lados?

\textbf{Resolução:}
A área é proporcional ao quadrado do lado.
\[
\sqrt{\frac{9}{16}} = \frac{3}{4}
\]

\textbf{Resposta:} $3:4$

\subsection*{Questão 13}
Se 8 torneiras enchem um tanque em 6 horas, em quanto tempo 12 torneiras iguais encherão o mesmo tanque?

\textbf{Resolução:}
Grandezas inversamente proporcionais.
\[
8 \cdot 6 = 12 \cdot x
\Rightarrow x = 4
\]

\textbf{Resposta:} 4 horas

\subsection*{Questão 14}
Quatro operários constroem um muro em 15 dias trabalhando 6 horas por dia. Quantos dias serão necessários para 6 operários trabalhando 5 horas por dia?

\textbf{Resolução:}
\[
4 \cdot 15 \cdot 6 = 6 \cdot x \cdot 5
\Rightarrow 360 = 30x
\Rightarrow x = 12
\]

\textbf{Resposta:} 12 dias

\subsection*{Questão 15}
Seis máquinas produzem 900 peças em 15 dias. Quantas máquinas são necessárias para produzir 1.200 peças em 10 dias?

\textbf{Resolução:}
\[
6 \cdot 15 \cdot 900 = x \cdot 10 \cdot 1200
\Rightarrow x = 6{,}75
\]

\textbf{Resposta:} 7 máquinas (aproximadamente)

\subsection*{Questão 16}
A razão entre as idades de Ana e Beatriz é $2:3$. Daqui a 4 anos, essa razão será $3:4$. Determine a idade atual de Ana.

\textbf{Resolução:}
\[
\frac{2x + 4}{3x + 4} = \frac{3}{4}
\Rightarrow 8x + 16 = 9x + 12
\Rightarrow x = 4
\]
Idade de Ana:
\[
2x = 8
\]

\textbf{Resposta:} 8 anos

\subsection*{Questão 17}
A razão entre as velocidades de dois atletas é $4:5$. Se o mais rápido percorre 100 m em 20 s, quanto tempo o outro leva?

\textbf{Resolução:}
\[
\frac{4}{5} = \frac{20}{x}
\Rightarrow 4x = 100
\Rightarrow x = 25
\]

\textbf{Resposta:} 25 s

\subsection*{Questão 18}
Três grandezas diretamente proporcionais $A$, $B$ e $C$ assumem os valores $A=4$, $B=6$ e $C=10$. Se $A$ passa a valer 10, determine o novo valor de $C$.

\textbf{Resolução:}
\[
\frac{4}{10} = \frac{10}{x}
\Rightarrow 4x = 100
\Rightarrow x = 25
\]

\textbf{Resposta:} 25

\subsection*{Questão 19}
Uma máquina A produz 120 peças em 6 horas, enquanto a máquina B produz a mesma quantidade em 4 horas. Determine a razão entre as eficiências das máquinas.

\textbf{Resolução:}
\[
A = \frac{120}{6} = 20 \quad B = \frac{120}{4} = 30
\]
\[
\frac{20}{30} = \frac{2}{3}
\]

\textbf{Resposta:} $2:3$

\subsection*{Questão 20}
Explique por que a regra de três é uma aplicação direta do conceito de proporção.

\textbf{Resposta:}
Porque se baseia na igualdade entre duas razões para determinar um valor desconhecido.

\end{multicols}

\end{document}
